\chapter{Ingénieur logiciel}\label{ch:in:software-engineer}

En février 2026, Virtu Financial, l'un des plus grands teneurs de marché électroniques cotés, comptait environ 1\,027 salariés répartis
dans treize pays. Cinq mois plus tard, Microsoft annonçait 77\,000 personnes rien que dans la recherche et le développement produit,
soixante-quinze fois l'effectif total de Virtu. Un ingénieur logiciel dans une société de trading est l'un de quelques centaines d'ingénieurs au plus,
à quelques pas du compte de résultat, et d'astreinte pendant que les marchés sont ouverts~; dans une entreprise technologique, le même ingénieur est
l'un de dizaines de milliers, loin du chiffre d'affaires, avec un produit qui fonctionne en permanence. Ce chapitre décrit les métiers
d'ingénierie dans les sociétés de trading, montre en quoi ils diffèrent des mêmes métiers dans la technologie, et compare les rémunérations que les dépôts
et l'enquête permettent d'établir pour eux.

\begin{center}\footnotesize
\begin{tabular}{@{}p{2.4cm}p{3.3cm}p{3.3cm}p{3.3cm}@{}}
\toprule
\multicolumn{4}{@{}l}{\textbf{Fiches de rôle~: trois spécialités d'ingénierie}}\\
\midrule
 & \omterm{def:in:software-engineer:lowlat}{ingénieur basse latence} & \omterm{def:in:software-engineer:fpga}{ingénieur FPGA} & \omterm{def:in:software-engineer:sre}{ingénieur de fiabilité des sites}\\
\midrule
travaille sur & la voie logicielle, des données de marché à l'ordre & la logique câblée sur matériel programmable sur cette voie & les systèmes de production et leur exploitation\\
unités & microsecondes et leurs queues & nanosecondes et cycles d'horloge & disponibilité, incidents, temps de rétablissement\\
rend compte à & responsable de la technologie de trading & responsable du matériel & responsable de l'infrastructure\\
codes dans les dépôts & 15-1252 & 15-1252, 17-2061 & 15-1252, 15-1299.08\\
enseigné dans & livre~13, ch.~1~; livre~14, ch.~1 & livre~14, ch.~6--7 & livre~15, ch.~28~; livre~16, ch.~21\\
\bottomrule
\end{tabular}
\end{center}

\section{Ingénierie cœur, basse latence, infrastructure, matériel et réseau}

Les ingénieurs d'une société de trading se répartissent selon ce qu'ils construisent. Les ingénieurs \emph{cœur} construisent la logique du système de trading et
les outils qui l'entourent~: le moteur de stratégies, la gestion des ordres, les contrôles de risque, la configuration que les traders modifient, et l'ingénierie
de la recherche du chapitre~19. Les autres se spécialisent selon la couche dont ils ont la charge.

\begin{definition}[Ingénieur basse latence]\label{def:in:software-engineer:lowlat}
Un \emph{ingénieur basse latence}\index{ingénieur basse latence} (low-latency engineer) est un ingénieur logiciel dont le travail consiste à raccourcir et à stabiliser le
temps écoulé entre un événement de marché et la réponse de la société à celui-ci~: le code, la disposition de la mémoire, les réglages du système d'exploitation
et la pile réseau sur ce chemin, mesurés en microsecondes et dans la queue de leur distribution.
\end{definition}

Le métier de l'\omterm{def:in:software-engineer:lowlat}{ingénieur basse latence} fait l'objet du livre~13~: d'où vient la latence (chapitre~1), comment les processeurs, les caches
et le système d'exploitation s'y ajoutent, et le contournement du noyau (kernel bypass) qui retire le noyau du chemin réseau (chapitre~13).
Le travail se mesure, il ne se discute pas~: une modification est mise en production lorsque son histogramme de latence, rejoué sur les mêmes données, est meilleur
là où cela compte, généralement dans la queue.

\begin{definition}[Ingénieur FPGA]\label{def:in:software-engineer:fpga}
Un \emph{ingénieur FPGA}\index{ingénieur FPGA} est un ingénieur qui conçoit, vérifie et maintient la logique destinée aux
circuits logiques programmables (FPGA)~: les feed handlers, les constructeurs d'ordres et les contrôles de risque qui
s'exécutent en matériel, où le temps écoulé entre l'arrivée d'un paquet et la réponse se compte en cycles d'horloge.
\end{definition}

Le livre~14 couvre le matériel (chapitres~6 et~7) et le réseau qui l'entoure~: commutateurs, cartes réseau, synchronisation
temporelle, colocation. L'\omterm{def:in:software-engineer:fpga}{ingénieur FPGA} écrit dans un langage de description de matériel, le vérifie en simulation bien plus à fond
que le logiciel n'est testé (un bogue matériel ne se corrige pas au milieu d'une journée de trading), et travaille avec
les ingénieurs réseau qui gèrent les câbles, les commutateurs et les baies colocalisées. C'est la plus petite des spécialités
d'ingénierie, et celle qui est la plus propre au trading.

Deux autres groupes complètent la carte. Les \emph{ingénieurs réseau} gèrent les liaisons entre la société et ses plateformes~: les baies colocalisées et les
interconnexions (cross-connect) (livre~14, chapitre~9), les commutateurs et les équipements de couche~1 (chapitre~2), les horloges qui horodatent
chaque paquet (chapitre~4), la capture qui les enregistre (chapitre~5), et, pour les sociétés qui se disputent la course entre villes, les
liaisons longue distance par fibre et par voie hertzienne (chapitres~13 et~14). Les \emph{ingénieurs d'infrastructure} gèrent ce sur quoi la recherche et le trading
tournent en dehors de la bourse~: les grappes de calcul (cluster) et leurs ordonnanceurs, le stockage des données tick, les systèmes de
construction et de mise en production, et les comptes cloud que certaines sociétés utilisent désormais pour la recherche (livre~15). Les deux groupes se mesurent à la capacité,
au coût et aux pannes plutôt qu'aux microsecondes, mais une panne dans l'un ou l'autre arrête le trading tout aussi sûrement qu'un bogue dans la stratégie.

\begin{definition}[Ingénieur de fiabilité des sites]\label{def:in:software-engineer:sre}
Un \emph{ingénieur de fiabilité des sites}\index{ingénieur de fiabilité des sites} (site reliability engineer) est un ingénieur logiciel qui assure la fiabilité des
systèmes de production~: la supervision et les alertes, la rotation d'astreinte, la gestion et la revue des incidents, la capacité, et l'automatisation qui
supprime le travail opérationnel manuel.
\end{definition}

Le nom vient de Google, dont le livre consacré à cette pratique la décrit comme «~ce qui se produit lorsqu'on demande à un ingénieur logiciel
de concevoir une équipe d'exploitation~»~; il plafonne le travail opérationnel à la moitié du temps de l'équipe, afin que l'autre moitié
serve à le supprimer par l'ingénierie. Les sociétés de trading appellent souvent ce métier ingénierie de production ou support de trading~;
les outils sont ceux du livre~15, chapitre~28~: les objectifs de niveau de service (SLO), la rotation d'astreinte et les revues sans recherche de coupable (blameless).

\section{En quoi une société de trading diffère d'une entreprise technologique}

Quatre différences façonnent la journée d'un ingénieur.
\begin{itemize}
\item \textbf{Distance à l'argent.} Le système d'une société de trading gagne ou perd directement de l'argent, dans les secondes qui suivent une modification.
Un ingénieur peut voir le résultat d'une mise en production dans le compte de résultat du lendemain, et une panne peut coûter davantage en quelques minutes
que le salaire annuel de l'ingénieur. L'ingénieur d'une entreprise technologique est en général à plusieurs couches du chiffre d'affaires.
\item \textbf{L'horloge du marché.} Un système de trading doit être disponible pendant que ses marchés sont ouverts et peut être modifié pendant qu'ils sont
fermés. Les mises en production partent entre deux séances, l'astreinte compte surtout lorsque les marchés sont ouverts, et les week-ends servent à la maintenance.
Les marchés de cryptoactifs, qui ne ferment jamais, suppriment ce rythme (chapitre~26).
\item \textbf{Taille.} Une société de trading de mille personnes compte quelques centaines d'ingénieurs~; les équipes produit et les équipes plateforme
du livre~16, chapitre~19, comptent quelques personnes chacune, et un seul ingénieur peut être responsable d'un système de bout en bout.
\item \textbf{Ce qui ne s'achète pas.} Une société de trading construit certains des éléments que d'autres entreprises achètent sous forme de service, lorsque la latence,
le contrôle ou le secret l'exige (livre~16, chapitre~20)~; et une grande partie de ce qu'elle construit, elle ne le publie pas.
\end{itemize}

\begin{example}[Un budget d'erreur sur l'horloge du marché]\label{ex:in:software-engineer:budget}
Un objectif de niveau de service de 99.9~\% de disponibilité mesuré sur toute l'année autorise 525.6 minutes de panne. Mesuré
sur les seules heures de négociation du marché actions américain (6.5 heures sur 252 jours, soit 1\,638 heures), le même objectif de 99.9~\% autorise 98.3 minutes, et 99.99~\%
autorise 9.8. La société de trading dispose d'un budget plus petit, et chacune de ses minutes tombe pendant que le marché bouge.
\end{example}

\section{Rémunération et progression comparées à celles des entreprises technologiques}

Les éléments sur les salaires proviennent de deux sources qui répondent à des questions différentes. Les demandes de certification de conditions de travail du
chapitre~14 comparent les sociétés de trading aux banques et aux bourses pour le même code de métier et le même \omterm{def:in:pay-levels-by-role-firm-type-and-seniority:soc}{niveau de salaire}, avec des
intervalles~; elles ne couvrent pas ici les employeurs technologiques, car le seul passage sur les dossiers n'a retenu que les employeurs financiers
documentés. L'enquête sur les professions compare les secteurs, y compris la technologie, mais seulement sous forme de centiles publiés, sans
intervalles.

\begin{dated}{2025-09}{Ce que gagnent les ingénieurs logiciels~: les données publiques}\label{dat:in:software-engineer:pay}
Demandes de certification de conditions de travail de l'exercice 2025 codées comme développeurs logiciels (15-1252), \omterm{def:in:how-pay-works:base}{salaire de base} médian proposé~: fonds systématiques
207\,000~\$ (67 demandes, 6 employeurs), plateformes 200\,000~\$ (35, 3), teneurs de marché 175\,000~\$ (153, 9), banques
155\,688~\$ (4\,790, 8), bourses 147\,900~\$ (296, 4). Sociétés de trading par rapport aux banques, rapport des médianes 1.19 (intervalle à 95~\%
1.16--1.23), 1.17 au niveau IV. Enquête sur les professions, mai 2025, médiane des développeurs logiciels~: portails de recherche sur internet et autres
services d'information 213\,520~\$~; éditeurs de logiciels 164\,550~\$~; secteur des valeurs mobilières 163\,290~\$~; infrastructure informatique et traitement
des données 156\,950~\$~; banques 136\,620~\$~; conception de systèmes informatiques 132\,050~\$. \omterm{def:in:how-pay-works:base}{Salaire de base} uniquement.
\end{dated}

\begin{omfigure}
\begin{tikzpicture}
\begin{axis}[width=10.5cm,height=5.2cm,xmin=80,xmax=320,ymin=0.4,ymax=5.6,ytick=data,
  yticklabels={{bourse},{banque},{teneur de marché},{plateforme},{fonds systématique}},yticklabel style={font=\footnotesize},
  xticklabel style={font=\footnotesize},xlabel={\small salaire de base proposé, en milliers de \$},grid=major,grid style={gray!20},
  table/col sep=comma]
\addplot[draw=none,mark=none,error bars/.cd,x dir=both,x explicit,error mark=none,error bar style={omDef,line width=1pt}]
  table[x=p50,y=pos,x error plus expr=\thisrow{p90}-\thisrow{p50},x error minus expr=\thisrow{p50}-\thisrow{p10}]
  {figdata/industry/20-software-engineer/kinds.csv};
\addplot[draw=none,mark=none,error bars/.cd,x dir=both,x explicit,error mark=none,error bar style={omDef!40,line width=7pt}]
  table[x=p50,y=pos,x error plus expr=\thisrow{p75}-\thisrow{p50},x error minus expr=\thisrow{p50}-\thisrow{p25}]
  {figdata/industry/20-software-engineer/kinds.csv};
\addplot[only marks,mark=|,mark size=5pt,line width=1.5pt,omThm] table[x=p50,y=pos]{figdata/industry/20-software-engineer/kinds.csv};
\end{axis}
\end{tikzpicture}

\omcaption{\omterm{def:in:how-pay-works:base}{Salaire de base} proposé aux développeurs logiciels dans les demandes de certification de conditions de travail de l'exercice 2025, par type d'employeur (code
15-1252, tous intitulés)~: 10\up{e} à 90\up{e}~centile (trait fin), intervalle interquartile (trait épais), médiane (repère). Données~:
\texttt{data/industry/lca\_soc.csv}, par l'intermédiaire de \texttt{in\_engineer.cells}.}\label{fig:in:software-engineer:kinds}
\end{omfigure}

La fourchette des teneurs de marché est la plus large~: leur 90\up{e}~centile, 295\,000~\$, est le plus élevé de tous les types, et leur 10\up{e}~centile
est proche de celui des banques. Regroupée et comparée niveau par niveau, la prime des sociétés de trading peut se résumer en un seul nombre, avec
un intervalle.

\begin{omfigure}
\begin{tikzpicture}
\begin{axis}[width=10.5cm,height=5.2cm,xmin=0.5,xmax=5.5,ymin=0.7,ymax=2.3,xtick={1,2,3,4,5},
  xticklabels={tous niveaux,I,II,III,IV},xlabel={\small niveau de salaire},ylabel={\small rapport des médianes},
  tick label style={font=\footnotesize},grid=major,grid style={gray!20},table/col sep=comma,unbounded coords=jump,
  legend style={legend cell align=left,font=\scriptsize,at={(0.5,-0.3)},anchor=north,/tikz/every even column/.append style={column sep=8pt}},
  legend columns=3]
\addplot[only marks,mark=o,black!60,thick,error bars/.cd,y dir=both,y explicit]
  table[x expr=\thisrow{pos}-0.15,y=bank2021,y error plus expr=\thisrow{bank2021_hi}-\thisrow{bank2021},
  y error minus expr=\thisrow{bank2021}-\thisrow{bank2021_lo}]{figdata/industry/20-software-engineer/ratio.csv};
\addplot[only marks,mark=*,omDef,thick,error bars/.cd,y dir=both,y explicit]
  table[x=pos,y=bank2025,y error plus expr=\thisrow{bank2025_hi}-\thisrow{bank2025},
  y error minus expr=\thisrow{bank2025}-\thisrow{bank2025_lo}]{figdata/industry/20-software-engineer/ratio.csv};
\addplot[only marks,mark=square*,omThm,thick,error bars/.cd,y dir=both,y explicit]
  table[x expr=\thisrow{pos}+0.15,y=exch2025,y error plus expr=\thisrow{exch2025_hi}-\thisrow{exch2025},
  y error minus expr=\thisrow{exch2025}-\thisrow{exch2025_lo}]{figdata/industry/20-software-engineer/ratio.csv};
\draw[black!50] (axis cs:0.5,1) -- (axis cs:5.5,1);
\legend{par rapport aux banques 2021,par rapport aux banques 2025,par rapport aux bourses 2025}
\end{axis}
\end{tikzpicture}

\omcaption{Médiane du \omterm{def:in:how-pay-works:base}{salaire de base} proposé dans les demandes de développeurs logiciels des sociétés de trading (teneurs de marché, fonds systématiques,
plateformes) rapportée à celle des banques et des bourses, globalement et par \omterm{def:in:pay-levels-by-role-firm-type-and-seniority:soc}{niveau de salaire}, avec des intervalles de confiance à 95~\% obtenus par bootstrap. Les points manquants sont
des comparaisons portant sur moins de dix demandes ou sur moins de trois employeurs d'un côté. Données~: \texttt{data/industry/lca\_kind\_gap.csv},
par l'intermédiaire de \texttt{in\_engineer.kind\_gaps}.}\label{fig:in:software-engineer:ratio}
\end{omfigure}

La prime des sociétés de trading par rapport aux banques représente environ un cinquième à chaque niveau où elle peut être mesurée pour l'exercice 2025, et elle
s'est réduite d'environ un tiers pour l'exercice 2021, les offres des banques ayant augmenté plus vite (la \cref{fig:in:software-engineer:ratio}). Par
rapport aux bourses, elle est plus grande. Le \omterm{def:in:how-pay-works:base}{salaire de base} sous-estime l'écart~: les ratios implicites entre \omterm{def:in:how-pay-works:base}{rémunération totale} et \omterm{def:in:how-pay-works:base}{salaire de base} du chapitre~14 étaient
plusieurs fois plus élevés chez les teneurs de marché que chez les banques.

\begin{omfigure}
\begin{tikzpicture}
\begin{axis}[width=10.5cm,height=5.6cm,xmin=0.7,xmax=5.3,ymin=60,ymax=300,xtick={1,2,3,4,5},
  xticklabels={10\up{e},25\up{e},médiane,75\up{e},90\up{e}},xlabel={\small centile},ylabel={\small salaire annuel, en milliers de \$},
  tick label style={font=\footnotesize},grid=major,grid style={gray!20},table/col sep=comma,
  legend style={font=\scriptsize,at={(0.5,-0.3)},anchor=north,legend cell align=left,/tikz/every even column/.append style={column sep=6pt}},
  legend columns=3]
\addplot[omThm,thick,mark=square*] table[x=pos,y=search]{figdata/industry/20-software-engineer/survey.csv};
\addplot[omDef,thick,mark=*] table[x=pos,y=securities]{figdata/industry/20-software-engineer/survey.csv};
\addplot[black!70,thick,mark=triangle*] table[x=pos,y=publishers]{figdata/industry/20-software-engineer/survey.csv};
\addplot[omDef!50,thick,dashed,mark=o,mark options={solid}] table[x=pos,y=banks]{figdata/industry/20-software-engineer/survey.csv};
\addplot[black!40,thick,dashed,mark=diamond*,mark options={solid}] table[x=pos,y=systems]{figdata/industry/20-software-engineer/survey.csv};
\legend{recherche sur internet et information,valeurs mobilières,éditeurs de logiciels,banques,conception de systèmes informatiques}
\end{axis}
\end{tikzpicture}

\omcaption{Salaires annuels des développeurs logiciels par secteur d'activité à cinq centiles, enquête sur les professions de mai 2025. Le secteur des valeurs mobilières
rémunère comme les éditeurs de logiciels à la médiane et au-dessus d'eux aux deux extrémités~; la recherche sur internet et l'information rémunère davantage
à chaque centile. Données~: \texttt{data/industry/oews\_roles.csv}, par l'intermédiaire de \texttt{in\_engineer.survey}.}\label{fig:in:software-engineer:survey}
\end{omfigure}

L'enquête place le secteur des valeurs mobilières à côté des éditeurs de logiciels, et non au-dessus d'eux~: le rapport de leurs centiles
vaut 1.27 au 10\up{e}, 0.99 à la médiane et 1.12 au 90\up{e} (la \cref{fig:in:software-engineer:survey}). Le secteur qui
regroupe les plus grandes entreprises d'internet rémunère davantage à chaque centile~: la médiane du secteur des valeurs mobilières vaut 0.77 de
la leur. Deux précautions s'imposent. Le secteur des valeurs mobilières, dans l'enquête, comprend les courtiers et les sociétés de gestion ainsi que les
sociétés de trading, dont les demandes se situent au-dessus de la médiane de l'enquête~; et l'enquête mesure des salaires, non les attributions d'actions qui constituent
une part importante de la rémunération d'un ingénieur de la technologie (chapitre~13). La conclusion honnête est étroite~: pour le \omterm{def:in:how-pay-works:base}{salaire de base}, une société de trading se mesure
à la partie la mieux rémunérée de l'industrie technologique, et non au-dessus d'elle.

La progression professionnelle diffère selon la taille. Un employeur qui compte des dizaines de milliers d'ingénieurs a besoin d'une échelle de niveaux pour les comparer
entre eux~; une société de quelques centaines de personnes peut juger chaque ingénieur sur les systèmes qu'il possède, et sa rémunération suit
l'année de la société (chapitre~13). Le chapitre~28 retrace les parcours.

\section{Ingénierie publique~: blogs, conférences et logiciel libre}

La plus grande part du code d'une société de trading est privée, mais certaines sociétés publient beaucoup, et ce qu'elles publient constitue la meilleure
preuve publique de ce que font leurs ingénieurs.

\begin{dated}{2026-09}{Ingénierie publique des sociétés de trading}\label{dat:in:software-engineer:public}
Jane Street écrit «~tout ce que nous pouvons en OCaml~», qu'elle appelle sa «~plateforme de développement principale~», publie
des bibliothèques et outils OCaml libres (Base, le système de construction Dune, les extensions du langage OxCaml, et Hardcaml, «~un
langage pour concevoir et simuler du matériel en OCaml~»), a financé des travaux sur le compilateur OCaml et son gestionnaire de paquets,
et tient un blog d'ingénierie public~; son organisation GitHub recensait 411 dépôts publics au 29 septembre 2026. Firedancer de Jump
Crypto, un client validateur pour la blockchain Solana, est libre~; son README indique que «~Le modèle de concurrence
s'inspire de l'expérience acquise dans le domaine du trading à basse latence~».
\end{dated}

Pour un candidat, l'ingénierie publique répond aux questions que les entretiens ne permettent pas de poser~: dans quel langage l'entreprise écrit, comment elle
teste, ce qu'elle considère comme un problème difficile. Pour l'entreprise qui recrute, c'est aussi un moyen de faire améliorer par des ingénieurs extérieurs les
outils dont elle dépend. Ce qui n'est jamais publié, c'est la logique de trading, et l'ingénieur qui rejoint la société accepte que l'essentiel de
son travail reste privé (le chapitre~28 examine ce qu'un candidat peut présenter à partir d'un tel poste).

\section{Tutoriel~: sociétés de trading, banques et technologie}\label{tut:in:software-engineer:tut}

\begin{tutorial}
\textbf{Objectif.} Comparer la rémunération des ingénieurs logiciel dans les sociétés de trading à celle des banques, des bourses et des secteurs technologiques, en tenant compte de
l'incertitude que permet chaque source.
\textbf{État final~:} \cref{fig:in:software-engineer:kinds,fig:in:software-engineer:ratio,fig:in:software-engineer:survey}.
\begin{tutsteps}
\item \textbf{Cellules par type.} Le fichier \texttt{lca\_soc.csv} du chapitre~19 pour le code 15-1252, par type d'employeur, à tous les niveaux.
\item \textbf{Rapport des médianes.} \texttt{in\_lca\_kind\_gap\_derive.py} regroupe les sociétés de trading et applique
\texttt{firm.roles.median\_gap} face aux banques et face aux bourses, globalement et par niveau, à partir du cache du chapitre~14,
en ne conservant une comparaison que lorsque chaque côté compte dix dépôts provenant de trois employeurs.
\item \textbf{Rapports de centiles.} \texttt{percentile\_ratio(a, b)} divise deux distributions publiées centile par
centile et ne renvoie rien là où une valeur est marquée (\cref{lst:in:software-engineer:ratio}).
\omcode{code/firm/roles/firm_roles.py}{284}{293}{Deux distributions de salaires publiées, comparées centile par centile.}\label{lst:in:software-engineer:ratio}
\item \textbf{Échelle.} Les effectifs tirés des dépôts donnent les rapports de l'accroche.
\end{tutsteps}
\end{tutorial}

Pour l'exercice 2025, les sociétés de trading par rapport aux banques~: 1.19 (1.16--1.23) globalement, 1.25 (1.14--1.38) au niveau~II, 1.18 (1.12--1.27) au
niveau~III, 1.17 (1.11--1.29) au niveau~IV~; l'intervalle du niveau~I (0.78--1.43) en contient un. Face aux bourses~: 1.25 globalement et
1.30 au niveau~IV. Dans l'enquête, le secteur des valeurs mobilières face aux éditeurs de logiciels~: 1.27, 1.07, 0.99, 1.00, 1.12 du 10\up{e}~centile au
90\up{e}~centile.

\textbf{Pour aller plus loin.} Ajouter de grands employeurs technologiques aux dépôts à l'aide d'une liste sourcée, ce qui demande une deuxième
passe sur les fichiers~; pondérer la comparaison par niveau pour que le rapport global ne soit pas un effet de composition (chapitre~19)~; ajouter les primes
et les actions tirées de la \omterm{def:in:pay-levels-by-role-firm-type-and-seniority:median}{rémunération du salarié médian} des déclarations de procuration (chapitre~14).

\section{Construction~: fiches d'ingénierie et comparaison par centiles}\label{bld:in:software-engineer:roles}

\begin{build}
\bfield{Rôle} Ajouter les trois spécialités d'ingénierie au registre des rôles et une comparaison des distributions publiées.
\bfield{Interface} \texttt{firm.roles}~: les fiches \texttt{low-latency engineer}, \texttt{FPGA engineer},
\texttt{site reliability engineer}~; \texttt{percentile\_ratio(a, b, keys) -> \{key: ratio or None\}}~; avec
\texttt{median\_gap} (chapitre~19) pour les données qui ont des valeurs.
\bfield{Règles} Un centile marqué ou manquant ne donne aucun rapport, jamais une supposition~; les rapports de centiles publiés
n'ont pas d'intervalle, et le chapitre le dit.
\bfield{Tests d'acceptation} \texttt{code/firm/roles/tests/}~: rapports connus~; une marque de plafonnement donne \texttt{None}.
\bfield{Pour aller plus loin} Bornes pour un rapport lorsqu'un côté est plafonné~; interpoler une distribution à partir de ses cinq centiles et
comparer les parts au-dessus d'un seuil.
\end{build}

\begin{omsources}
\item Virtu Financial, Form 10-K for 2025~; Flow Traders, Annual Report 2025~; Microsoft, Form 10-K for fiscal 2026.
\item Beyer, B., C.~Jones, J.~Petoff et N.~R.~Murphy (dir., 2016), \emph{Site Reliability Engineering}, O'Reilly.
\item Page technologie de Jane Street et organisation GitHub~; dépôt Firedancer.
\item Tableaux du chapitre~14 tirés des fichiers LCA du Department of Labor~; enquête du BLS sur les professions, mai 2025.
\end{omsources}

\section{Exercices}

\begin{exercise}[$\star$]\label{exo:in:software-engineer:1}
Calculer combien de fois l'effectif de Virtu représente l'effectif de recherche et développement produit de Microsoft, et combien de fois il représente son effectif total.
\end{exercise}

\begin{exercise}[$\star$]\label{exo:in:software-engineer:2}
Calculer combien de minutes de panne par an un objectif de 99.95~\% autorise sur les heures de négociation des actions américaines.
\end{exercise}

\begin{exercise}[$\star$]\label{exo:in:software-engineer:3}
Le rapport annuel 2025 de Flow Traders donne son effectif de fin 2024 à 609 ETP dans une section et à 625 dans une autre. Dire ce que doit faire
un lecteur d'un tel écart.
\end{exercise}

\begin{exercise}[$\star\star$]\label{exo:in:software-engineer:4}
Expliquer pourquoi la prime des sociétés de trading sur les bourses est plus élevée au niveau~II qu'au niveau~IV.
\end{exercise}

\begin{exercise}[$\star\star$]\label{exo:in:software-engineer:5}
Le rapport entre la médiane des sociétés de trading et celle des banques est passé de 1.32 pour l'exercice 2021 à 1.19 pour l'exercice 2025. En utilisant les médianes du
tableau dérivé (165\,000~\$ et 125\,000~\$ en 2021~; 185\,000~\$ et 155\,688~\$ en 2025), dire quel côté a le plus évolué.
\end{exercise}

\begin{exercise}[$\star\star$]\label{exo:in:software-engineer:6}
Un ingénieur qui hésite entre une société de trading et une grande entreprise technologique ne compare que les \omterm{def:in:how-pay-works:base}{salaires de base}. Dire ce que cette comparaison
laisse de côté de chaque côté.
\end{exercise}

\begin{exercise}[$\star\star\star$]\label{exo:in:software-engineer:7}
\emph{Programmation.} Avec \texttt{percentile\_ratio}, comparer les développeurs logiciel du secteur des valeurs mobilières à chacun des
cinq autres secteurs du tableau de l'enquête. Dire pour combien des 25 comparaisons par centile le secteur des valeurs mobilières
verse une rémunération plus élevée.
\end{exercise}

\begin{exercise}[$\star\star\star$]\label{exo:in:software-engineer:8}
\emph{Trouver la faille.} «~L'enquête montre que le secteur des valeurs mobilières ne paie pas ses développeurs logiciel plus que ne le font les
éditeurs de logiciels, donc les sociétés de trading rémunèrent leurs ingénieurs comme le secteur technologique.~»
\end{exercise}

\section{Problème~: société de trading ou entreprise technologique~?}

\begin{problem}[{Problème du week-end --- société de trading ou entreprise technologique~?}]\label{pb:in:software-engineer:1}
Un ingénieur logiciel qui compte cinq ans dans une grande entreprise technologique reçoit une offre d'un teneur de marché. Elle veut savoir en quoi
le travail et la rémunération diffèrent.

\medskip\noindent\textbf{Partie I --- Les métiers.}
\begin{enumerate}
\item Définir l'\omterm{def:in:software-engineer:lowlat}{ingénieur basse latence}, l'\omterm{def:in:software-engineer:fpga}{ingénieur FPGA} et l'\omterm{def:in:software-engineer:sre}{ingénieur de fiabilité des sites} (SRE).
\item Que construisent les core engineers~?
\item D'où vient le nom de site reliability engineer, et quelle règle sur le travail opérationnel l'accompagne~?
\item Qu'est-ce qui rend les tests de l'\omterm{def:in:software-engineer:fpga}{ingénieur FPGA} plus stricts que ceux d'un ingénieur logiciel~?
\item Quels livres de la série enseignent chaque spécialité~?
\end{enumerate}

\medskip\noindent\textbf{Partie II --- L'entreprise.}
\begin{enumerate}[resume]
\item Comparer la taille des deux employeurs aux chiffres de l'accroche.
\item Énoncer les quatre différences entre une société de trading et une entreprise technologique.
\item Calculer le budget d'erreur d'un objectif de 99.9~\% sur une année, puis sur les heures de négociation.
\item Quel effet l'horloge du marché a-t-elle sur les mises en production et sur l'astreinte~?
\item Que publient les sociétés de trading sur leur ingénierie, et que ne publient-elles pas~?
\end{enumerate}

\medskip\noindent\textbf{Partie III --- La rémunération.}
\begin{enumerate}[resume]
\item Que peut comparer chaque source (dépôts réglementaires, enquête), et que ne peut-elle pas comparer~?
\item Donner les médianes des sociétés de trading par catégorie, ainsi que celles des banques.
\item Donner les ratios par rapport aux banques par niveau, avec les intervalles.
\item Comparer l'industrie financière aux éditeurs de logiciels et à la recherche sur le web dans l'enquête.
\item Que laisse de côté le \omterm{def:in:how-pay-works:base}{salaire de base} de chaque côté~?
\end{enumerate}

\medskip\noindent\textbf{Partie IV --- Le verdict.}
\begin{enumerate}[resume]
\item Énoncer le \emph{résultat nommé}~: le ratio entre le \omterm{def:in:how-pay-works:base}{salaire de base} médian proposé par les sociétés de trading et celui des banques pour les développeurs logiciels,
à chaque \omterm{def:in:pay-levels-by-role-firm-type-and-seniority:soc}{niveau de salaire}, avec les intervalles bootstrap, ainsi que la comparaison de l'enquête avec les secteurs technologiques.
\item Pourquoi la comparaison avec les banques est-elle plus nette que la comparaison avec la technologie~?
\item Comment diffère la progression~?
\item Que doit-elle demander au teneur de marché au sujet de l'astreinte~?
\item En deux phrases, répondez-lui.
\end{enumerate}
\end{problem}

\section{Questions d'entretien}

\begin{interviewq}[$\star$ \iqroles{développeur}]\label{iq:in:software-engineer:1}
Qu'est-ce que la latence de queue, et pourquoi une société de trading se soucie-t-elle davantage du 99.9\up{e}~centile que de la moyenne~?
\end{interviewq}

\begin{interviewq}[$\star$ \iqroles{développeur}]\label{iq:in:software-engineer:2}
Un service doit être disponible 99.9~\% des heures de négociation. Quelle est la durée d'indisponibilité tolérée sur une année, et comment l'utiliseriez-vous~?
\end{interviewq}

\begin{interviewq}[$\star\star$ \iqroles{développeur}]\label{iq:in:software-engineer:3}
Comment le contournement du noyau réduit-il la latence, et à quoi renonce-t-on~?
\end{interviewq}

\begin{interviewq}[$\star\star$ \iqroles{développeur, trader}]\label{iq:in:software-engineer:4}
Votre système a envoyé des ordres en double pendant deux secondes à l'ouverture. Décrivez la première heure qui suit.
\end{interviewq}

\begin{interviewq}[$\star\star$ \iqroles{développeur}]\label{iq:in:software-engineer:5}
Quand placeriez-vous un contrôle de risque dans un FPGA plutôt que dans le logiciel~?
\end{interviewq}

\begin{interviewq}[$\star\star\star$ \iqroles{développeur}]\label{iq:in:software-engineer:6}
Concevez le processus de mise en production d'un système de trading qui ne doit pas changer pendant que le marché est ouvert, mais qui doit être livré chaque jour.
\end{interviewq}
